In recent years, the capabilities of artificial intelligence models have grown rapidly, along with the complexity of their architectures. 
This progress has enabled the solution of tasks that once appeared unattainable, such as the use of language models to solve complex mathematical problems \parencite{frontierMath}. 
At the same time, current systems have reached scales far beyond those of a few years ago, both in terms of the number of parameters and structural complexity \parencite{EpochLargeScaleModels2024}. 
This lack of transparency is especially problematic in sensitive settings, such as medical diagnosis or high-stakes decision-making, where producing an accurate prediction is not enough: it is equally important to understand how that prediction was made. In response to this challenge, the field of Explainable Artificial Intelligence (XAI) has emerged \parencite{mersha2024explainable}, aiming to develop methods capable of making model decisions more interpretable.

Within the field of explainability, one of the most influential approaches is feature attribution, where each input variable is assigned a score reflecting its contribution to the prediction. 
In this line of work, SHAP is one of the most well-known methods \parencite{shapOriginalPaper}. It is based on the Shapley values, a concept originating in cooperative game theory \parencite{shapley1953value}. 
Intuitively, Shapley values provide a fair way to distribute the total value achieved by a coalition among its participants, according to their average marginal contribution when joining. 
Translated to machine learning, features play the role of players, and the value assigned to each feature aims to quantify its contribution to the model’s prediction for a given instance.

Despite its popularity, SHAP has important limitations. On the one hand, depending on the family of ML models under consideration and the underlying distribution of the data, the computation of the Shapley values can be tractable or not. For example,
it is well-known that for product distributions the complexity of computing these scores is equivalent to the complexity of computing the average 
value of the model \parencite{van2022tractability}. In particular, this implies that there exist efficient algorithms for the case where the models are given, for example, by decision trees or by restricted families of circuits \parencite{arenas2023complexity}, while the problem is intractable whenever the models are more expressive, as in the case of neural nets. 
For more general families of distributions the problem becomes much harder: for instance, when considering Naive Bayes distributions, the computation of this score is intractable already when considering trivial models that only depend on one feature \parencite{van2022tractability}.

On the other hand, there are also concerns about the quality of the explanations it produces. In particular, it is not always clear how to interpret, in the context of AI models, the axioms inherited from game theory, and several works point out that these explanations may fail to adequately capture relationships 
such as dependence, relevance, or causality between variables \parencite{fryer2021shapley, marques2023logic, huang2023inadequacy}. More generally, there is no absolute consensus on what should be considered an “explanation” in AI: many current techniques are pragmatic approximations to a much broader problem related to understanding, causality, and human interpretability \parencite{MILLER20191, lipton2017mythosmodelinterpretability}.

The Asymmetric Shapley Values (ASV from now on) are a proposal introduced by \cite{frye2019asymmetric} to overcome some of these limitations. This variant incorporates causal information into the computation of attribution scores by restricting attention to permutations of variables that respect a causal graph. 
In this way, instead of treating all features symmetrically, priority is given to the more fine-grained ones (in the sense that the information they provide is not present in the other variables). This can produce explanations that are more faithful to the data-generating process and can help better detect phenomena such as bias or discrimination. 
Moreover, although ASV was primarily proposed to improve explanatory quality, it naturally raises the question of whether this additional restriction may also provide computational benefits.

Therefore, in this paper we address the problem of efficiently computing the ASV scores, which was not explored thoroughly in previous literature. We will prove that for some families of models for which
computing the Shapley values is intractable there are still polynomial time algorithms for computing the ASV. Inspired by this initial result, we develop an algorithm to compute the ASV exactly in more general contexts
by exploiting the fact that many permutations of the features can be grouped in a unique equivalence class to simplify computations. Finally, we also propose a flexible approximation scheme that allows to approximate the ASV as long as it is possible to (1) Sample a topological order of the causal graph uniformly at random, and (2) Sample a data point with respect to the data distribution after conditioning a subset of the features to have specific values. We instantiate all these general algorithms with particular 
examples of causal graphs and machine learning models and show experimentally that they are efficient.
